\documentclass[12pt,a4paper]{article}

\usepackage[T1]{fontenc}
\usepackage[utf8]{inputenc}
\usepackage{lmodern}
\usepackage[left=30mm,right=25mm,top=25mm,bottom=25mm]{geometry}
\usepackage[english, ngerman]{babel}
\usepackage[autostyle=true,german=quotes]{csquotes}
\usepackage[onehalfspacing]{setspace}
\usepackage[expansion=false]{microtype}
\usepackage{booktabs}
\usepackage{tabularx}
\usepackage{enumitem}
\usepackage{xcolor}
\usepackage{amsmath}
\usepackage{amssymb}
\usepackage{pifont}
\usepackage{etoolbox}
\usepackage{caption}
\usepackage{tikz}
\usetikzlibrary{positioning,patterns,arrows.meta}
\usepackage{pgfplots}
\pgfplotsset{compat=1.18}
% Diagrammfarben (geprüft auf Farbsehschwächen), zusätzlich Marker/Strichart
\definecolor{cblue}{HTML}{2A78D6}
\definecolor{corange}{HTML}{EB6834}
\definecolor{ctext}{HTML}{52514E}
\usepackage[backend=biber,style=ieee,urldate=long]{biblatex}
\addbibresource{sycl.bib}
\usepackage{hyperref}

\hypersetup{
    colorlinks,
    citecolor   = black,
    filecolor   = black,
    linkcolor   = black,
    urlcolor    = black,
    pdftitle    = {SYCL: Portable heterogene Programmierung mit C++},
    pdfauthor   = {Erlind Sejdiu},
    pdfcreator  = {pdflatex},
    pdfproducer = {LaTeX with hyperref}
}

% Tabellen, Abbildungen und Literaturverzeichnis einzeilig
\AtBeginEnvironment{tabular}{\singlespacing}
\AtBeginEnvironment{tabularx}{\singlespacing}
\AtBeginEnvironment{tikzpicture}{\singlespacing}
\captionsetup{font=small, labelfont=bf}
\AtBeginBibliography{\singlespacing}

% Markierung für Stellen, die noch ergänzt/geprüft werden müssen
\newcommand{\offen}[1]{\textcolor{red}{[#1]}}

\begin{document}

% ---------------------------------------------------------------------------
% Titelei (römische Seitenzahlen, Titelblatt ohne Nummer)
% ---------------------------------------------------------------------------
\pagenumbering{Roman}

\begin{titlepage}
    \raggedright
    Erlind Sejdiu\\
    106357
    \vspace*{3cm}

    \begin{center}
        {\LARGE\bfseries SYCL\\[0.3em]
        Portable heterogene Programmierung mit C++\\[0.3em]
        \large Problemstellung, Anwendungsfälle und Programmiermodell\par}
        \vspace{2cm}
        {\large Seminararbeit\\
        der Hochschule Karlsruhe\\
        im Master Studiengang Informatik\\
        im Wintersemester 2026/27\par}
    \end{center}
    \vfill
    \begin{tabular}{@{}ll}
        Betreuer: & \offen{Name der betreuenden Person}\\
        Themengruppe: & \offen{Themengruppe des Seminars}\\
        Abgabedatum: & \offen{Datum}
    \end{tabular}
\end{titlepage}

\setcounter{page}{1}
\tableofcontents
\newpage
\listoffigures
\listoftables
\newpage

% ---------------------------------------------------------------------------
% Hauptteil (arabische Seitenzahlen ab 1)
% ---------------------------------------------------------------------------
\pagenumbering{arabic}

\section{Einleitung}

Rechenleistung wächst heute vor allem durch Spezialisierung: Neben der CPU
stehen in fast jedem Rechner GPUs und andere Beschleuniger, die datenparallele
Aufgaben um Größenordnungen schneller bearbeiten. Die Software hält mit dieser
Vielfalt jedoch nur schwer Schritt, denn jeder Hersteller bringt sein eigenes
Programmiermodell mit. SYCL ist ein offener Standard der Khronos Group, der
dieses Problem lösen will: Ein einziger Quellcode in Standard-C++ soll auf CPUs
und GPUs verschiedener Hersteller laufen \cite{khronos2020sycl}.

Diese Arbeit erklärt, welches Problem SYCL adressiert, für welche
Anwendungsfälle es gedacht ist und wie sein Programmiermodell diese erfüllt.
Ein kurzes Beispiel macht die Konzepte greifbar.

% ---------------------------------------------------------------------------
\section{Problemstellung}\label{sec:problem}
% ---------------------------------------------------------------------------

Heutige Systeme kombinieren CPUs mit GPUs unterschiedlicher Hersteller. Für
die Software ergeben sich daraus drei Probleme:

\begin{enumerate}[nosep]
    \item \textbf{Herstellerbindung:} CUDA läuft nur auf NVIDIA-GPUs
    \cite{nickolls2008cuda}; AMD (HIP) und Intel (Level~Zero) haben eigene
    Schnittstellen.
    \item \textbf{Mehrfache Codebasen:} Wer mehrere Hersteller unterstützen
    will, pflegt mehrere Implementierungen parallel
    (Abbildung~\ref{fig:codebasen}, links).
    \item \textbf{Abstraktion gegen Kontrolle:} Offene Alternativen sind
    entweder sehr hardwarenah und aufwendig wie OpenCL
    \cite{stone2010opencl} oder lassen wie OpenMP-Direktiven wenig Kontrolle
    über Speicher und Ausführung.
\end{enumerate}

\begin{figure}[htbp]
\centering
\begin{tikzpicture}[
    code/.style={draw, rounded corners=2pt, minimum width=2.3cm, minimum height=0.55cm, font=\scriptsize, fill=gray!10},
    hw/.style={draw, minimum width=1.25cm, minimum height=0.6cm, align=center, font=\scriptsize, fill=corange!15, draw=corange},
    >={Stealth[length=1.8mm]}
]
    \node[font=\footnotesize\bfseries] at (1.9,2.75) {Herstellerspezifisch};
    \node[code] (c1) at (0.9,2) {CUDA-Code};
    \node[code] (c2) at (0.9,1) {HIP-Code};
    \node[code] (c3) at (0.9,0) {OpenMP-Code};
    \node[hw, minimum width=1.6cm] (h1) at (3.6,2) {NVIDIA-GPU};
    \node[hw, minimum width=1.6cm] (h2) at (3.6,1) {AMD-GPU};
    \node[hw, minimum width=1.6cm] (h3) at (3.6,0) {CPU};
    \draw[->] (c1) -- (h1); \draw[->] (c2) -- (h2); \draw[->] (c3) -- (h3);
    \draw[dashed, gray] (5.1,-0.4) -- (5.1,2.9);
    \node[font=\footnotesize\bfseries] at (9.3,2.75) {Mit SYCL};
    \node[code, fill=cblue!15, draw=cblue, minimum width=4.2cm] (s) at (9.3,2) {ein SYCL-Quellcode (C++17)};
    \node[code, minimum width=4.2cm] (i) at (9.3,1) {DPC++ \;|\; AdaptiveCpp};
    \node[hw] (g1) at (6.7,0) {NVIDIA};
    \node[hw] (g2) at (8.43,0) {AMD};
    \node[hw] (g3) at (10.17,0) {Intel};
    \node[hw] (g4) at (11.9,0) {CPU};
    \draw[->] (s) -- (i);
    \foreach \g in {g1,g2,g3,g4} \draw[->] (i.south) -- (\g.north);
\end{tikzpicture}
\caption{Mehrere herstellerspezifische Codebasen (links) gegenüber einer SYCL-Codebasis, die über Backends wie CUDA, HIP, Level~Zero oder OpenCL auf verschiedener Hardware läuft (rechts)}
\label{fig:codebasen}
\end{figure}

% ---------------------------------------------------------------------------
\section{SYCL als Lösungsansatz}\label{sec:sycl}
% ---------------------------------------------------------------------------

SYCL ist ein offener Standard der Khronos Group auf Basis von C++17
\cite{khronos2020sycl}. Implementierungen wie DPC++ \cite{intel_llvm} und
AdaptiveCpp \cite{alpay2020hipsycl, adaptivecpp} bilden ihn über Backends wie
CUDA, HIP, Level~Zero, OpenCL oder OpenMP auf die Hardware ab und übersetzen
den Gerätecode meist erst beim ersten Start \emph{just in time} für das
konkrete Gerät \cite{alpay2023sscp}. Den drei Problemen setzt SYCL vier
Entwurfsziele entgegen:

\begin{description}[style=sameline, leftmargin=1.5em, nosep]
    \item[Portabilität:] Ein Quellcode, viele Backends; das Gerät wird erst
    beim Übersetzen bzw.\ zur Laufzeit gewählt.
    \item[Offenheit:] Ein offener Standard mit mehreren unabhängigen
    Implementierungen statt eines Herstellerprodukts.
    \item[Produktivität:] \emph{Single-Source} in Standard-C++: Host- und
    Gerätecode stehen in einer Datei, Kernel sind gewöhnliche Lambdas, und
    Templates und Typen werden geteilt.
    \item[Kontrolle:] Wo nötig, legt das Modell Hardwarekonzepte wie
    Work-Groups und lokalen Speicher offen, um gezielt optimieren zu können.
\end{description}

\noindent Tabelle~\ref{tab:vergleich} vergleicht SYCL entlang dieser Ziele mit
den Alternativen. Kokkos \cite{edwards2014kokkos} und RAJA
\cite{beckingsale2019raja} verfolgen ähnliche Ziele, sind aber kein offener
Industriestandard.

\begin{table}[htbp]
\centering
\small
\newcommand{\ja}{\ding{51}}
\newcommand{\nein}{--}
\newcommand{\teils}{($\checkmark$)}
\begin{tabular}{@{}lcccc@{}}
\toprule
 & \textbf{herstellerunabh.} & \textbf{offener Standard} & \textbf{Standard-C++} & \textbf{Kontrolle} \\
\midrule
CUDA            & \nein   & \nein   & \teils & \ja \\
OpenCL          & \ja     & \ja     & \nein  & \ja \\
OpenMP-Offload  & \ja     & \ja     & \ja    & \teils \\
Kokkos/RAJA     & \ja     & \nein   & \ja    & \teils \\
\textbf{SYCL}   & \ja     & \ja     & \ja    & \ja \\
\bottomrule
\end{tabular}
\caption{SYCL im Vergleich zu anderen Programmiermodellen (($\checkmark$) = eingeschränkt)}
\label{tab:vergleich}
\end{table}

% ---------------------------------------------------------------------------
\section{Anwendungsfälle}\label{sec:usecases}
% ---------------------------------------------------------------------------

Aus den Entwurfszielen ergeben sich vier typische Einsatzfelder
(Tabelle~\ref{tab:usecases}): wissenschaftliches Rechnen auf Systemen
verschiedener Hersteller, etwa auf dem Supercomputer Aurora mit Intel-GPUs
oder in der Molekulardynamik-Software GROMACS
\cite{deakin2020sycl, breyer2022svm}; die Migration bestehender
CUDA-Anwendungen, unterstützt durch das Werkzeug SYCLomatic; die gemeinsame
Nutzung von CPU und GPU in einem Rechner; und portable Bibliotheken wie
oneMKL oder oneDPL \cite{reinders2023dpcpp}. Die Tabelle zeigt, welche
Anforderung jedes Feld stellt und welcher Teil des Programmiermodells sie
erfüllt.

\begin{table}[htbp]
\centering
\footnotesize
\begin{tabularx}{\textwidth}{@{}>{\raggedright\arraybackslash}p{2.6cm}>{\raggedright\arraybackslash}X>{\raggedright\arraybackslash}X@{}}
\toprule
\textbf{Anwendungsfall} & \textbf{Anforderung} & \textbf{Antwort im Programmiermodell} \\
\midrule
Rechnen über Herstellergrenzen & ein Code für viele Systeme, hohe Leistung & Backends und Geräteauswahl (\ref{sec:geraete}); hardwarenahe Hierarchie (\ref{sec:ausfuehrung}); Interoperabilität mit Herstellerbibliotheken (\ref{sec:erweitert}) \\
Migration von CUDA & vertraute Konzepte, schrittweiser Umbau & Work-Group/Sub-Group analog zu Block/Warp (\ref{sec:ausfuehrung}); zeigerbasiertes USM (\ref{sec:speicher}) \\
CPU und GPU gemeinsam & Überlappung, Ausweichgerät & asynchrone Queues und Abhängigkeitsgraph (\ref{sec:queues}); CPU als Gerät (\ref{sec:geraete}) \\
Portable Bibliotheken & generischer Code, Anpassung je Gerät & Single-Source mit Templates (\ref{sec:sycl}); \texttt{get\_info}, Aspekte, Spezialisierungskonstanten (\ref{sec:geraete}, \ref{sec:erweitert}) \\
\bottomrule
\end{tabularx}
\caption{Anwendungsfälle, ihre Anforderungen und die Mechanismen, die sie erfüllen}
\label{tab:usecases}
\end{table}

% ---------------------------------------------------------------------------
\section{Das Programmiermodell}\label{sec:modell}
% ---------------------------------------------------------------------------

\subsection{Plattformen, Geräte und Geräteauswahl}\label{sec:geraete}

SYCL unterscheidet den \emph{Host}, auf dem das Hauptprogramm läuft, und
\emph{Geräte}, auf denen Kernel ausgeführt werden. Geräte sind nach
\emph{Plattformen} gruppiert, wobei jede Plattform einem Backend entspricht;
auch die CPU kann als Gerät auftreten. Welches Gerät verwendet wird, bestimmt
ein \emph{Device Selector}. Vordefiniert sind Selektoren für das als am
leistungsfähigsten eingeschätzte Gerät, für GPUs und für CPUs; eigene
Selektoren vergeben Punktwerte anhand von Geräteeigenschaften. Über
\emph{Aspekte} kann ein Programm zudem prüfen, ob ein Gerät optionale
Fähigkeiten wie doppelte Genauigkeit oder bestimmte USM-Arten unterstützt,
und mit \texttt{get\_info} Eigenschaften wie die maximale Work-Group-Größe
abfragen \cite{khronos2020sycl}. Implementierungen erlauben außerdem eine
Auswahl von außen, etwa über Umgebungsvariablen, sodass derselbe Code ohne
Neuübersetzung auf einem anderen Gerät läuft. Dass auch die CPU als
vollwertiges Gerät gilt, macht sie zum Ausweichgerät, wenn kein
Beschleuniger vorhanden ist, und erlaubt die Entwicklung ohne GPU.
Abbildung~\ref{fig:selektor} veranschaulicht die Auswahl.

\begin{figure}[htbp]
\centering
\begin{tikzpicture}[
    dev/.style={draw, rounded corners=2pt, minimum width=3.9cm, minimum height=0.55cm, font=\scriptsize, fill=gray!10, anchor=west},
    score/.style={font=\scriptsize, text=ctext, anchor=west},
    >={Stealth[length=1.8mm]}
]
    \node[font=\footnotesize\bfseries, anchor=west] at (0,2.2) {verfügbare Geräte (Backend)};
    \node[dev, fill=cblue!15, draw=cblue, very thick] (d1) at (0,1.5) {NVIDIA-GPU (CUDA)};
    \node[dev] (d2) at (0,0.75) {Intel-GPU (Level Zero)};
    \node[dev] (d3) at (0,0) {CPU (OpenCL / OpenMP)};
    \node[score] at (4.05,1.5) {Wertung 900};
    \node[score] at (4.05,0.75) {Wertung 700};
    \node[score] at (4.05,0) {Wertung 300};
    \node[draw, rounded corners=2pt, align=center, font=\scriptsize, minimum height=1.2cm, fill=white] (sel) at (8.2,0.75) {Device Selector\\[1pt]bewertet Eigenschaften\\und Aspekte};
    \node[draw, rounded corners=2pt, font=\scriptsize, fill=cblue!15, draw=cblue, minimum height=0.6cm] (q) at (12.0,0.75) {Queue};
    \draw[->] (5.95,1.5) -- (sel.west |- 0,1.05);
    \draw[->] (5.95,0.75) -- (sel.west);
    \draw[->] (5.95,0) -- (sel.west |- 0,0.45);
    \draw[->, thick, cblue] (sel) -- node[above, font=\scriptsize, text=black, align=center] {höchste\\Wertung} (q);
\end{tikzpicture}
\caption{Geräteauswahl: Der Selector bewertet jedes Gerät, die Queue wird an das bestbewertete gebunden}
\label{fig:selektor}
\end{figure}

\subsection{Queues, Command Groups und Abhängigkeiten}\label{sec:queues}

Die Verbindung zu einem Gerät stellt eine \emph{Queue} her. Arbeit wird in
Form von \emph{Command Groups} eingereicht: Eine Command Group beschreibt eine
Aktion, etwa einen Kernel oder eine Kopie, zusammen mit den Daten, die sie
benötigt. Technisch ist eine Command Group eine Funktion, die
ein \emph{Handler}-Objekt erhält: Über den Handler deklariert sie ihre
Datenzugriffe und Abhängigkeiten und startet genau eine Aktion. Das Einreichen
kehrt sofort zurück und liefert ein \emph{Event}. Die Ausführung ist also
asynchron, und der Host wartet nur, wenn er es ausdrücklich verlangt.
Abhängigkeiten zu früheren Aufträgen lassen sich über deren Events angeben.
Mit \emph{Host-Tasks} kann zudem gewöhnlicher Host-Code als Knoten in den
Graphen eingereiht werden, etwa um nach einer Berechnung eine Datei zu
schreiben, ohne den Rest des Programms zu blockieren.

Aus den Datenzugriffen und Events baut die Laufzeitumgebung einen gerichteten
azyklischen Abhängigkeitsgraphen auf (Abbildung~\ref{fig:dag}). Voneinander
unabhängige Aufträge darf sie parallel oder in anderer Reihenfolge ausführen,
abhängige startet sie erst, wenn ihre Vorgänger fertig sind
\cite{reinders2023dpcpp}. So lassen sich Datentransfers, Berechnungen und
Host-Arbeit überlappen, ohne Threads oder Streams von Hand zu verwalten. Wer
eine strikte Reihenfolge bevorzugt, legt die Queue als \emph{in-order} an.

\begin{figure}[htbp]
\centering
\begin{tikzpicture}[
    node/.style={draw, rounded corners=2pt, minimum height=0.6cm, inner sep=3pt, font=\scriptsize, fill=gray!10},
    kern/.style={node, fill=cblue!15, draw=cblue},
    >={Stealth[length=2mm]}
]
    \node[node] (ca) at (0,0.5) {Kopie $A$ $\rightarrow$ Gerät};
    \node[node] (cb) at (0,-0.5) {Kopie $B$ $\rightarrow$ Gerät};
    \node[kern] (k1) at (3.7,0) {Kernel 1: $C = A \cdot B$};
    \node[node] (cd) at (3.7,-1.2) {Kopie $D$ $\rightarrow$ Gerät};
    \node[kern] (k2) at (7.4,-0.6) {Kernel 2: $E = C + D$};
    \node[node] (ce) at (10.9,-0.6) {Kopie $E$ $\rightarrow$ Host};
    \draw[->] (ca) -- (k1);
    \draw[->] (cb) -- (k1);
    \draw[->] (k1) -- (k2);
    \draw[->] (cd) -- (k2);
    \draw[->] (k2) -- (ce);
\end{tikzpicture}
\caption{Abhängigkeitsgraph: Die Kopie von $D$ kann parallel zu Kernel~1 laufen}
\label{fig:dag}
\end{figure}

\subsection{Ausführungsmodell}\label{sec:ausfuehrung}

Ein datenparalleler Kernel wird mit \texttt{parallel\_for} für jeden Punkt
eines ein- bis dreidimensionalen Indexraums einmal ausgeführt; eine einzelne
Ausführung heißt \emph{Work-Item}. Bei einer einfachen \texttt{range} gibt der
Programmierer nur die Gesamtgröße an und überlässt die Aufteilung der
Implementierung. Bei einer \texttt{nd\_range} legt er zusätzlich fest, wie
die Work-Items zu \emph{Work-Groups} gebündelt werden
(Abbildung~\ref{fig:ndrange}). Die Work-Items einer Gruppe laufen gemeinsam
auf einer Recheneinheit, teilen sich schnellen lokalen Speicher und können
sich über Barrieren synchronisieren. Seit SYCL~2020 gibt es darunter
\emph{Sub-Groups}, deren Work-Items im Gleichschritt laufen und Daten direkt
austauschen können \cite{khronos2020sycl}.

\begin{figure}[htbp]
\centering
\begin{tikzpicture}[scale=0.75, font=\footnotesize]
    \fill[cblue!15] (1,2) rectangle (2,3);
    \draw[step=0.25, gray!60] (1,2) grid (2,3);
    \fill[cblue] (1.25,2.5) rectangle (1.5,2.75);
    \draw[step=1, gray] (0,0) grid (4,4);
    \draw[thick] (0,0) rectangle (4,4);
    \draw[thick, cblue] (1,2) rectangle (2,3);
    \draw[->] (4.4,3.4) node[right, align=left] {Work-Group: teilt lokalen Speicher,\\synchronisiert über Barrieren} -- (2.05,2.85);
    \draw[->] (4.4,1.6) node[right, align=left] {Work-Item: eine Ausführung\\des Kernels} -- (1.45,2.55);
    \node[below] at (2,0) {Indexraum der \texttt{nd\_range}};
\end{tikzpicture}
\caption{Gliederung des Indexraums in Work-Groups und Work-Items}
\label{fig:ndrange}
\end{figure}

\begin{figure}[htbp]
\centering
\begin{tikzpicture}[
    sy/.style={draw, rounded corners=2pt, minimum width=2.8cm, minimum height=0.55cm, font=\scriptsize, fill=cblue!15, draw=cblue},
    hw/.style={draw, rounded corners=2pt, minimum width=3.6cm, minimum height=0.55cm, font=\scriptsize, fill=gray!10},
]
    \node[font=\footnotesize\bfseries] at (0,3.0) {SYCL};
    \node[font=\footnotesize\bfseries] at (4.6,3.0) {NVIDIA-GPU};
    \node[font=\footnotesize\bfseries] at (9.2,3.0) {CPU};
    \foreach \y/\s/\g/\c in {2.3/nd\_range/Grid (gesamte GPU)/alle Kerne,
                              1.5/Work-Group/Thread-Block auf einem SM/ein Kern (Thread),
                              0.7/Sub-Group/Warp (32 Threads)/SIMD-Vektor,
                              -0.1/Work-Item/Thread/eine SIMD-Lane} {
        \node[sy] (s) at (0,\y) {\s};
        \node[hw] (g) at (4.6,\y) {\g};
        \node[hw] (c) at (9.2,\y) {\c};
        \draw[dashed, gray] (s) -- (g);
        \draw[dashed, gray] (g) -- (c);
    }
\end{tikzpicture}
\caption{Typische Abbildung der SYCL-Ausführungshierarchie auf GPU und CPU}
\label{fig:abbildung}
\end{figure}

Diese Hierarchie bildet die Hardware ab (Abbildung~\ref{fig:abbildung}) und
entspricht weitgehend dem Modell von CUDA \cite{nickolls2008cuda}: Auf einer NVIDIA-GPU wird eine Work-Group zu
einem Thread-Block mit höchstens 1024 Threads \cite{nvidia_cuda_guide}, eine
Sub-Group zu einem Warp. Auf einer CPU arbeitet typischerweise ein Kern eine
Work-Group ab, und die Sub-Group entspricht den Lanes der Vektoreinheit.
Neben \texttt{parallel\_for} bietet SYCL~2020 fertige Gruppenalgorithmen, etwa
um innerhalb einer Work-Group einen Wert zu verteilen oder zu summieren, sowie
\emph{Reduktionen}, mit denen sich Ergebnisse aller Work-Items, etwa eine
Summe oder ein Maximum, ohne eigene Synchronisation zusammenfassen lassen.

Für den Gerätecode gelten Einschränkungen, weil er auf Hardware ohne
vollständige C++-Laufzeitumgebung laufen muss: Ausnahmen, virtuelle
Funktionsaufrufe, Rekursion und dynamische Speicherreservierung sind im Kernel
nicht erlaubt, und Daten werden per Wertkopie übernommen
\cite{khronos2020sycl}.

\subsection{Speichermodell}\label{sec:speicher}

GPUs besitzen meist eigenen Speicher, der vom Hauptspeicher des Hosts getrennt
ist. SYCL bietet zwei Modelle, um Daten zwischen beiden zu bewegen
(Tabelle~\ref{tab:speicher}).

\begin{table}[htbp]
\centering
\small
\begin{tabularx}{\textwidth}{@{}l>{\raggedright\arraybackslash}X>{\raggedright\arraybackslash}X@{}}
\toprule
 & \textbf{Buffer und Accessoren} & \textbf{Unified Shared Memory (USM)} \\
\midrule
Zugriff        & über Accessor-Objekte mit Zugriffsmodus & über gewöhnliche Zeiger \\
Datentransfer  & automatisch durch die Laufzeit & \emph{shared}: automatisch; \emph{device}: explizit mit \texttt{memcpy} \\
Abhängigkeiten & automatisch aus den Accessoren & selbst über Events oder in-order-Queue \\
Stärke         & sicher, wenig Fehlerquellen & vertraut aus C++, leicht zu portieren \\
\bottomrule
\end{tabularx}
\caption{Die beiden Speichermodelle von SYCL}
\label{tab:speicher}
\end{table}

\begin{figure}[htbp]
\centering
\begin{tikzpicture}[
    box/.style={draw, rounded corners=2pt, minimum height=0.6cm, inner sep=3pt, font=\scriptsize, fill=gray!10},
    kern/.style={box, fill=cblue!15, draw=cblue},
    lab/.style={font=\scriptsize, align=center},
    lz/.style={font=\scriptsize\itshape, text=corange, align=center},
    >={Stealth[length=1.8mm]}
]
    \node[box] (h) at (0,1.2) {Host-Daten};
    \node[box] (b) at (2.3,1.2) {Buffer $X$};
    \node[kern] (k1) at (6.5,1.2) {Kernel 1: schreibt $X$};
    \node[kern] (k2) at (11.3,1.2) {Kernel 2: liest $X$};
    \node[box] (ha) at (11.3,-0.4) {Host-Accessor};
    \node[box] (h2) at (6.5,-0.4) {Ergebnis auf dem Host};
    \draw[->] (h) -- (b);
    \draw[->] (b) -- node[lab, above] {Accessor\\(write)} node[lz, below] {Kopie zum\\Gerät} (k1);
    \draw[->] (k1) -- node[lab, above] {Accessor\\(read)} node[lz, below] {Abhängigkeit} (k2);
    \draw[->] (k2) -- node[lz, right] {Kopie zurück} (ha);
    \draw[->] (ha) -- (h2);
\end{tikzpicture}
\caption{Buffer-Modell: Aus den Accessoren leitet die Laufzeit Kopien und Abhängigkeiten selbst ab (orange)}
\label{fig:buffer}
\end{figure}

\begin{figure}[htbp]
\centering
\begin{tikzpicture}[font=\scriptsize, >={Stealth[length=2mm]}]
    \draw[fill=gray!8, rounded corners=2pt] (0,0) rectangle (3.2,3.2);
    \node[anchor=north, font=\footnotesize\bfseries] at (1.6,3.15) {Host (CPU)};
    \draw[fill=white] (0.3,0.3) rectangle (2.9,2.5);
    \node[align=center] at (1.6,1.4) {Hauptspeicher\\[2pt]USM \emph{host}};
    \draw[fill=gray!8, rounded corners=2pt] (6.0,0) rectangle (12.6,3.2);
    \node[anchor=north, font=\footnotesize\bfseries] at (9.3,3.15) {Gerät (GPU)};
    \draw[fill=white] (6.3,0.3) rectangle (12.3,2.5);
    \node[anchor=north] at (9.3,2.48) {globaler Speicher (USM \emph{device})};
    \foreach \x in {6.55,9.4} {
        \draw[fill=cblue!15, draw=cblue] (\x,0.5) rectangle (\x+2.65,1.95);
        \node at (\x+1.325,1.7) {lokaler Speicher};
        \node at (\x+1.325,1.4) {(Work-Group)};
        \foreach \i in {0,1,2,3}
            \draw[fill=cblue, draw=cblue] (\x+0.3+\i*0.55,0.65) rectangle (\x+0.7+\i*0.55,0.95);
    }
    \node[below] at (9.3,0) {dunkel: privater Speicher je Work-Item};
    \draw[<->, thick] (2.9,1.9) -- node[above] {\texttt{memcpy}} (6.3,1.9);
    \draw[<->, thick, dashed, corange] (2.9,0.9) -- node[below, align=center, text=black] {USM \emph{shared}:\\automatische Migration} (6.3,0.9);
\end{tikzpicture}
\caption{Speicherhierarchie von Host und Gerät und die Wege des Datentransfers}
\label{fig:speicher}
\end{figure}

Beim \emph{Buffer-Modell} kapselt ein Buffer die Daten, und Kernel greifen über
\emph{Accessoren} darauf zu, die angeben, ob gelesen oder geschrieben wird.
Daraus leitet die Laufzeit sowohl die nötigen Kopien als auch die
Abhängigkeiten im Graphen ab. Die Daten gehören während der Lebensdauer des
Buffers der Laufzeit; erst wenn der Buffer zerstört wird oder der Host über
einen \emph{Host-Accessor} zugreift, sind die Ergebnisse wieder im
Hauptspeicher sichtbar (Abbildung~\ref{fig:buffer}). Dieses Modell ist sicher,
verlangt aber ein Umdenken gegenüber gewöhnlichen Zeigern. Beim zeigerbasierten \emph{USM}, das mit
SYCL~2020 eingeführt wurde, wird Speicher ähnlich wie mit \texttt{malloc}
angelegt. Es gibt drei Arten: \emph{device} (nur auf dem Gerät),
\emph{host} (im Hauptspeicher, vom Gerät aus lesbar) und \emph{shared}
(auf beiden Seiten gültig, die Laufzeit migriert die Daten bei Bedarf)
\cite{reinders2023dpcpp}. USM erleichtert die Portierung bestehender
Programme, überträgt aber die Verantwortung für Abhängigkeiten auf den
Programmierer.

Auf dem Gerät unterscheidet SYCL zudem (Abbildung~\ref{fig:speicher})
\emph{globalen} Speicher (für alle
Work-Items sichtbar, groß, aber langsam), \emph{lokalen} Speicher (pro
Work-Group, klein und schnell) und \emph{privaten} Speicher (Register). Diese
Hierarchie gezielt zu nutzen ist, wie bei CUDA, der wichtigste Hebel für
Leistung \cite{hwu2022pmpp}.

\subsection{Fehlerbehandlung, Spezialisierung und Interoperabilität}\label{sec:erweitert}

Fehler meldet SYCL als C++-Ausnahmen. Synchrone Fehler, etwa eine ungültige
Work-Group-Größe beim Einreichen, werden sofort geworfen. Asynchrone Fehler,
die erst während der Ausführung auf dem Gerät auftreten, sammelt die
Laufzeit und übergibt sie einem \emph{Async Handler}, sobald der Host
wartet. Diese Trennung ist eine Folge der asynchronen Ausführung: Ein Fehler
in einem Kernel kann erst sichtbar werden, wenn der Host längst weitere
Aufträge eingereicht hat (Abbildung~\ref{fig:fehler}). \emph{Spezialisierungskonstanten} erlauben es, Werte wie eine
Kachelgröße erst zur Laufzeit festzulegen und dennoch vom Compiler als
Konstante behandeln zu lassen, was zur JIT-Kompilierung passt. Über die
\emph{Backend-Interoperabilität} lassen sich schließlich native Objekte des
Backends abrufen, um etwa herstellerspezifische Bibliotheken direkt auf
SYCL-Daten aufzurufen \cite{khronos2020sycl}.

\begin{figure}[htbp]
\centering
\begin{tikzpicture}[
    ev/.style={draw, rounded corners=2pt, minimum height=0.55cm, inner sep=3pt, font=\scriptsize, fill=gray!10, align=center},
    lz/.style={font=\scriptsize\itshape, text=corange, align=center},
    >={Stealth[length=1.8mm]}
]
    \node[font=\footnotesize\bfseries, anchor=east] at (-0.2,1.2) {Host};
    \node[font=\footnotesize\bfseries, anchor=east] at (-0.2,0) {Gerät};
    \draw[->, gray] (0,1.2) -- (13,1.2) node[right, font=\scriptsize, text=ctext] {Zeit};
    \draw[->, gray] (0,0) -- (13,0);
    \node[ev] (s1) at (1.1,1.2) {Kernel 1\\einreichen};
    \node[ev, fill=cblue!15, draw=cblue, minimum width=3.6cm] (k1) at (4.4,0) {Kernel 1 läuft \;\textcolor{corange}{\ding{55}} Fehler};
    \node[ev, fill=corange!15, draw=corange] (s2) at (4.4,1.2) {Kernel 2 einreichen\\(ungültige Work-Group)};
    \node[ev] (h) at (7.9,1.2) {weitere\\Host-Arbeit};
    \node[ev] (w) at (11.0,1.2) {\texttt{wait\_and\_throw()}\\$\rightarrow$ Async Handler};
    \draw[->] (s1.south) -- (2.6,0.28);
    \node[lz, above=1pt of s2] {synchron: Ausnahme sofort};
    \draw[->, dashed, corange, thick] (k1.east) to[out=0, in=-90] node[lz, below right, pos=0.45] {asynchron: erst beim Warten} (w.south);
\end{tikzpicture}
\caption{Synchrone Fehler werden beim Einreichen geworfen, asynchrone erst, wenn der Host wartet}
\label{fig:fehler}
\end{figure}

% ---------------------------------------------------------------------------
\section{Anwendungsbeispiel: Matrixmultiplikation}\label{sec:beispiel}
% ---------------------------------------------------------------------------

Zur Veranschaulichung wurde ein sequentielles C++-Programm für die
Matrixmultiplikation $C = A \cdot B$ \cite{hwu2022pmpp} schrittweise nach
SYCL portiert und
mit DPC++ auf einer CPU (AMD Ryzen~5 5600X) und einer GPU (NVIDIA GeForce
RTX~5060~Ti) ausgeführt. Der Weg entspricht dem Anwendungsfall der Migration:
Zuerst wurden die äußeren Schleifen durch ein \texttt{parallel\_for} über
USM \emph{shared} ersetzt, wobei die Rechenvorschrift unverändert zum Kernel
wurde. Danach wurden die Daten mit USM \emph{device} explizit kopiert, und
schließlich lud jede Work-Group Blöcke der Matrizen in den lokalen Speicher.

Alle Stufen liefen ohne Codeänderung auf beiden Geräten; gewechselt wurde nur
das Gerät beim Programmstart. Abbildung~\ref{fig:ergebnisse} zeigt zwei
Beobachtungen. Erstens kostet bei kleinen Matrizen der Datentransfer mehr als
die Rechnung, weshalb das Überlappen über den Abhängigkeitsgraphen wichtig
ist (a). Zweitens ist die beste Work-Group-Größe geräteabhängig: Auf der GPU
ist eine Kachelgröße von $T = 16$ am schnellsten, auf der CPU $T = 32$, und
bei $T = 64$ überschreitet die Work-Group die Grenze von 1024 Threads der
NVIDIA-GPU, sodass der Kernel dort nicht startet (b). Die Kennzahl für
Performance-Portabilität nach Pennycook et~al.\ \cite{pennycook2016metric}
fällt dann auf null.

\begin{figure}[htbp]
\centering
\pgfplotsset{
    every axis/.append style={
        font=\scriptsize, width=7cm, height=4.8cm,
        axis line style={ctext}, tick style={ctext},
        grid style={gray!20}, ymajorgrids,
        legend style={draw=none, fill=none, font=\scriptsize},
        legend cell align=left,
    }
}
\begin{minipage}[t]{0.49\textwidth}\centering
\begin{tikzpicture}
\begin{axis}[
    xmode=log, ymode=log, log basis x=2,
    xtick={256,512,1024,2048}, xticklabels={256,512,1024,2048},
    xlabel={Matrixgröße $n$}, ylabel={Zeit in ms},
    legend pos=north west,
    title={(a) Rechnung und Transfer, GPU},
]
\addplot[cblue, line width=1.2pt, mark=*, mark size=2pt]
    coordinates {(256,0.038) (512,0.263) (1024,1.974) (2048,15.869)};
\addlegendentry{Rechnung}
\addplot[corange, line width=1.2pt, dashed, mark=square*, mark size=2pt, mark options={solid}]
    coordinates {(256,0.173) (512,0.405) (1024,1.344) (2048,3.934)};
\addlegendentry{Transfer}
\end{axis}
\end{tikzpicture}
\end{minipage}\hfill
\begin{minipage}[t]{0.49\textwidth}\centering
\begin{tikzpicture}
\begin{axis}[
    ybar, bar width=7pt, enlarge x limits=0.18,
    ymin=0, ymax=115, ytick={0,25,50,75,100},
    symbolic x coords={8,16,32,64}, xtick=data,
    xticklabels={{$T{=}8$\\PP 0{,}59},{$T{=}16$\\PP 0{,}93},{$T{=}32$\\PP 0{,}99},{$T{=}64$\\PP 0}},
    xticklabel style={align=center},
    ylabel={Effizienz in \%},
    legend style={at={(0.5,1.0)}, anchor=south, legend columns=2},
    title style={yshift=3ex},
    title={(b) Effizienz je Kachelgröße, $n = 2048$},
]
\addplot[fill=cblue, draw=white, line width=0.8pt]
    coordinates {(8,83) (16,100) (32,99) (64,0)};
\addlegendentry{GPU}
\addplot[fill=corange, draw=white, line width=0.8pt, postaction={pattern=north east lines, pattern color=white}]
    coordinates {(8,45) (16,88) (32,100) (64,91)};
\addlegendentry{CPU}
\node[font=\scriptsize, text=ctext, rotate=90, anchor=west] at ([xshift=-4.5pt]axis cs:64,2) {startet nicht};
\end{axis}
\end{tikzpicture}
\end{minipage}
\caption{Datentransfer (a) und Leistung je Kachelgröße mit Performance-Portabilität (b)}
\label{fig:ergebnisse}
\end{figure}

% ---------------------------------------------------------------------------
\section{Bewertung und Fazit}\label{sec:fazit}
% ---------------------------------------------------------------------------

SYCL löst das Problem der Herstellerbindung und der mehrfachen Codebasen
überzeugend: Ein Quellcode in Standard-C++ läuft auf CPUs und GPUs
verschiedener Hersteller, und mehrere unabhängige Implementierungen sichern
die Offenheit. Die Anwendungsfälle tragen dabei verschiedene Teile des
Programmiermodells. Für die Migration aus CUDA sind die enge
Verwandtschaft der Ausführungshierarchie und das zeigerbasierte USM
entscheidend, für die gemeinsame Nutzung von CPU und GPU die asynchronen
Queues mit ihrem Abhängigkeitsgraphen und für portable Bibliotheken
Single-Source mit Templates sowie die Geräteabfragen.

Die Grenze liegt im dritten Problem, dem Zielkonflikt zwischen Abstraktion und
Kontrolle. SYCL abstrahiert die Schnittstelle, nicht die Hardware: Für hohe
Leistung muss der Programmierer Work-Groups und lokalen Speicher ebenso
verstehen wie bei CUDA, und Parameter wie die Work-Group-Größe bleiben
gerätespezifisch, wie das Beispiel zeigt. Hinzu kommen eine merkliche Abhängigkeit von der
Implementierung \cite{deakin2020sycl} und eine anspruchsvolle
Werkzeugkette.

Insgesamt erfüllt SYCL sein Versprechen \enquote{einmal schreiben, überall
ausführen} für die Korrektheit. Für hohe Leistung auf allen Geräten stellt es
mit Geräteabfragen und Spezialisierungskonstanten die Werkzeuge bereit, nimmt
dem Programmierer die gerätespezifische Abstimmung aber nicht ab.

\newpage
\printbibliography[heading=bibintoc, title={Literaturverzeichnis}]

\end{document}
